\chapter{اللقاء الثامن والثلاثون: أمثال الإنفاق والإخلاص وبيان ما يقبل في الصدقة}

\section{مقدمة}
السلام عليكم ورحمة الله وبركاته. الحمد لله رب العالمين، وأشهد أن لا إله إلا الله وحده لا شريك له، وأشهد أن محمداً عبده ورسوله، صلى الله عليه وعلى آله وصحبه وسلم.

استأنف الشيخ هذا اللقاء من قوله تعالى: \begin{ayah}مَثَلُ الَّذِينَ يُنْفِقُونَ أَمْوَالَهُمْ فِي سَبِيلِ اللَّهِ\end{ayah}، ودار المجلس على آيات النفقة والصدقة من جهة المعنى، والنية، والمنّ والأذى، والإخلاص، وأمثال القرآن في هذه الأبواب، حتى ختم بقوله تعالى: \begin{ayah}وَاعْلَمُوا أَنَّ اللَّهَ غَنِيٌّ حَمِيدٌ\end{ayah}.

وفي هذا اللقاء تتجلى مهارة \rawi{الطبري} في وصل السياق بعضه ببعض، ورد الآيات إلى أمها الأولى، ثم تفصيل الوجوه التي يبطل بها العمل أو يزكو، مع تحرير دقيق للفرق بين الواجب والتطوع في باب الصدقة.

\separator

\section{رد الآيات إلى أصلها في سياق السورة}
\isnad{قال أبو جعفر:} هذه الآية وما بعدها مردودة عند \rawi{الطبري} إلى قوله تعالى: \begin{ayah}مَنْ ذَا الَّذِي يُقْرِضُ اللَّهَ قَرْضًا حَسَنًا\end{ayah}، وأن ما توسط بينها من قصص بني إسرائيل، ومحاجة إبراهيم، والمرور على القرية، وسؤال إبراهيم ربه أن يريه كيف يحيي الموتى، جاء اعتراضاً إلهياً بديعاً لحكم متعددة:
\begin{itemize}
    \item إقامة الحجة على المشركين المكذبين بالبعث.
    \item تحريض المؤمنين على الجهاد والإنفاق والثقة بنصر الله.
    \item قطع عذر اليهود بما أطلع الله نبيه عليه من خفي أخبارهم.
    \item إعذار المنافقين وتحذيرهم من مصير من سبقهم.
\end{itemize}

ونبّه الشيخ إلى أن هذا من أنفس مواضع الكتاب في بيان صلة القصص بسياق السورة، وأن القصص في القرآن ليست وحدات منفصلة، بل ترد إلى مقاصد السورة التي جاءت فيها.

\begin{fawaid}[من أجمل وجوه التفسير رد المقاطع إلى أصلها في السياق]
قد تتخلل بين الآية وأختها قصص واعتراضات تبدو متفرقة، لكن المفسر البصير يرد الجميع إلى مقصد واحد يجمعها، وبهذا تنتظم السورة في القلب ولا تبدو كأنها موضوعات متناثرة.
\end{fawaid}

\section{مثل الذين ينفقون أموالهم في سبيل الله}
فسر \rawi{الطبري} قوله تعالى: \begin{ayah}مَثَلُ الَّذِينَ يُنْفِقُونَ أَمْوَالَهُمْ فِي سَبِيلِ اللَّهِ\end{ayah} بالنفقة على أنفسهم في جهاد أعداء الله، أو معونة المجاهدين في سبيله، وأن المثل المضروب هنا يراد به بيان عظم البركة والثواب، كما تنبت الحبة سبع سنابل في كل سنبلة مائة حبة.

ثم ذكر الخلاف في قوله: \begin{ayah}وَاللَّهُ يُضَاعِفُ لِمَنْ يَشَاءُ\end{ayah}، هل هذه الزيادة بعد السبعمائة تشمل غير المنفق في سبيل الله أو هي راجعة إلى المنفق نفسه؟ ورجح أن السياق في المنفقين في سبيل الله خاصة، وأن المضاعفة فوق السبعمائة لهم هم، لا لغيرهم.

\begin{fawaid}[السياق من أقوى ما يضبط مدلول المضاعفة والوعود]
إذا جاء وعد عظيم في باب مخصوص، فالأصل حمل ما بعده عليه ما لم يأت دليل يصرفه. ولهذا رد الطبري قوله: (والله يضاعف لمن يشاء) إلى المنفق في سبيل الله نفسه.
\end{fawaid}

\section{المنّ والأذى يفسدان كمال النفقة}
عند قوله تعالى: \begin{ayah}ثُمَّ لَا يُتْبِعُونَ مَا أَنْفَقُوا مَنًّا وَلَا أَذًى\end{ayah}، شرح \rawi{الطبري} أن المن هو إظهار الصنيعة والتفضل على المعطَى، والأذى هو شكايته أو تعييره أو ذكر عطائه له بما يؤلمه ويكسره.

ووقف الشيخ عند أمثلة السلف في هذا الباب، ومنها من كان يؤذي الفقير أو المجاهد قبل أن يعطيه، أو يذكر عطاياه بما يفسد نفعها على قلبه. ومن أعظم ما يربى عليه المؤمن هنا أن العطاء إذا لم يسلم من المن والأذى انقلب في كثير من أحواله وبالاً على صاحبه.

\begin{fawaid}[ليس كل عطاء قربة حتى يسلم من المن والأذى]
قد يبذل الإنسان مالاً كثيراً، ثم يفسد أثره أو أجره بتعظيم نفسه، أو تعيير من أعطاه، أو طلب المقابل المعنوي من الناس، فيخرج بذلك من روح الصدقة وإن بقيت صورتها.
\end{fawaid}

\section{معنى (لا خوف عليهم ولا هم يحزنون)}
بيّن الشيخ أن نفي الخوف والحزن عن هؤلاء لا يعني نفي أصل الخوف المحمود الذي يكون في الدنيا، فإن المؤمنين يخافون الله ويخافون يوماً عظيماً، لكن المراد هنا نفي الخوف الذي يضرهم عند لقاء الله، ونفي الحزن الذي يلحق أهل الحسرة على ما فاتهم.

فخوف الدنيا إذا حمل على العمل والطاعة نافع، أما الخوف يوم القيامة على المصير، والحزن على ضياع العمل، فهذا الذي أمنه الله من أنفق ماله له على الوجه المرضي.

\section{قول معروف ومغفرة خير من صدقة يتبعها أذى}
توقف الشيخ عند هذه الآية وقفة تربوية لطيفة، وبيّن أن الشرع لا يريد من المؤمن مجرد صورة الإحسان المالي، بل يريد اجتماع الإحسان القولي والقلبي معه. فقد يكون الاعتذار الجميل، أو الرد بلطف، أو التغافل عن تقصير المحتاج، خيراً من صدقة يعقبها كسر قلب أو إذلال نفس.

وفي ختام هذا الموضع جاء اسم الله \begin{ayah}غَنِيٌّ حَلِيمٌ\end{ayah}، فنبّه الشيخ إلى أن \rawi{الطبري} من أكثر الناس عناية بربط ختام الآية بمضمونها؛ فالله غني عن هذه الصدقات، حليم لا يعاجل المسيء بالعقوبة مع ما بدر منه في باب القربات.

\begin{fawaid}[خواتيم الآيات تربي القلوب على معاني الأحكام]
إذا ختم الله حكماً باسم من أسمائه أو فعل من أفعاله فليس ذلك فاصلة مجردة، بل هو مفتاح لفهم المراد وتربية القلب على الوجه الذي ينبغي أن يتلقى به الحكم.
\end{fawaid}

\section{المثلان في الرياء: الصفوان والجنة المحترقة}
من أجمل ما في هذا اللقاء تحليل \rawi{الطبري} للمثلين اللذين ضربهما الله للمنفق رياء الناس:
\begin{itemize}
    \item المثل الأول: \begin{ayah}كَمَثَلِ صَفْوَانٍ عَلَيْهِ تُرَابٌ\end{ayah}.
    \item المثل الثاني: \begin{ayah}أَيَوَدُّ أَحَدُكُمْ أَنْ تَكُونَ لَهُ جَنَّةٌ\end{ayah}.
\end{itemize}

ورجح \rawi{الطبري} أن المثل الثاني أيضاً راجع إلى المرائي، لا إلى مطلق المفرط في الطاعة، لأن الله ساق الآيات في هذا الباب متتابعة في بيان بطلان النفقة بالمن والرياء، فحمل المثل على هذا المعنى أولى من نقله إلى معنى آخر لم يجر له ذكر قريب.

وبيّن الشيخ دقة هذا الترجيح: فالمرائي له في الدنيا هيئة عمل ومنفعة عاجلة وثناء من الناس، كصاحب الجنة التي ينتفع بها زمناً، لكنه إذا جاء وقت حاجته الحقيقية إلى عمله، وجده قد احترق وذهب، كما احترقت الجنة على صاحبها عند الكبر وضعف الذرية.

\begin{fawaid}[قد يكون للباطل نفع عاجل لكنه ينهار عند شدة الحاجة]
من أعظم معاني هذا المثل أن الرياء قد يعطي صاحبه صورة نفع في الدنيا، ومدحاً، ووجاهة، لكنه عند ساعة الفقر الحقيقي إلى العمل لا يجد منه شيئاً، وتلك من أوجع الحسرات.
\end{fawaid}

\section{الطبري يرد بعض المعاني الصحيحة إذا لم تكن هي المراد بالسياق}
ذكر \rawi{الطبري} أقوالاً حسنة في تفسير مثل الجنة المحترقة، كحمله على من فرط في الطاعة ثم ختم له بغيرها، أو على من أفسد عمله بعد طول سير، ثم بيّن أن هذه المعاني في نفسها صحيحة، لكن حمل الآية عليها هنا بعيد عن السياق.

وهذا من نفاسة صنيعه؛ فإنه لا يرد كل معنى لأنه باطل، بل قد يرد المعنى مع الاعتراف بحسنه إذا رأى أن الآية لم تسق له في هذا الموضع.

\begin{fawaid}[ليس كل معنى صحيح يصلح تفسيراً للآية في موضعها]
قد يكون القول حسناً في نفسه، وله وجه معتبر، لكن المفسر المحقق ينظر: هل سيق له هذا الموضع أم لا؟ فصحة المعنى وحدها لا تكفي إن خالفه السياق.
\end{fawaid}

\separator

\section{كذلك يبين الله لكم الآيات لعلكم تتفكرون}
عدَّ الشيخ هذا الموضع من أجمل تعليقات \rawi{الطبري} في السورة؛ لأن الله لما أمر بالإنفاق لم يكتف بمجرد الأمر به، بل بيّن وجوهه، وشروطه، ومبطلاته، ومقاصده، وأثر النية فيه، وما يقبل منه وما لا يقبل.

ومن هنا استخرج الشيخ أصلاً عملياً عظيماً: أن العبد مأمور بأمرين متلازمين:
\begin{itemize}
    \item أن يطلب حكم الله ويتحراه ويسأل عنه.
    \item وأن يجتهد بعد العلم في العمل به.
\end{itemize}

فليس الورع أن يتشدد الإنسان بلا علم، ولا أن يعمل أعمالاً كثيرة وهو يجهل شروطها وموانعها، بل الورع أن يعرف ما أمر الله به ثم يقف عنده.

\begin{fawaid}[من تمام التقوى أن يبين الله لك ما تتقيه]
لا يمكن أن يتقي العبد الله في باب من الأبواب حتى يعرف حدوده وأحكامه، ولذلك كان من رحمة الله أن يبين تفاصيل الأوامر والنواهي، لا أن يترك الناس لأهوائهم وتخميناتهم.
\end{fawaid}

\section{يا أيها الذين آمنوا أنفقوا من طيبات ما كسبتم}
فسر \rawi{الطبري} قوله تعالى: \begin{ayah}مِنْ طَيِّبَاتِ مَا كَسَبْتُمْ\end{ayah} بالجياد من الأموال الحلال: من الذهب والفضة ونحو ذلك، ثم ألحق به قوله: \begin{ayah}وَمِمَّا أَخْرَجْنَا لَكُمْ مِنَ الْأَرْضِ\end{ayah} في الحب والثمار وسائر ما تجب فيه الزكاة.

ثم جاء النهي: \begin{ayah}وَلَا تَيَمَّمُوا الْخَبِيثَ مِنْهُ تُنْفِقُونَ\end{ayah}. وبيّن الشيخ أن الخبيث هنا ليس المراد به مجرد الشيء القديم أو المستعمل الذي لا يزال نافعاً، وإنما المقصود تعمد الرديء الذي لا يرضاه الإنسان لنفسه، أو أن يجعل معظم صدقته من هذا الذي لا خير فيه.

\begin{fawaid}[المعتبر في الخبيث هنا الرداءة المقصودة لا مجرد القدم]
ليس كل ما استعمله الإنسان أو استغنى عنه يدخل في الخبيث المنهي عنه، بل المدار على كونه مما لا يرضاه لنفسه لو كان هو الآخذ، أو مما لا نفع فيه، أو مما قصد به التملص من حق الله والعباد.
\end{fawaid}

\section{سبب النزول والقياس عليه}
ذكر \rawi{الطبري} سبب النزول في تعليق بعض الأنصار قنوان التمر الرديء والحشف في موضع الصدقة، فأنزل الله النهي عن تعمد هذا الرديء في الإنفاق الواجب.

ثم نبّه الشيخ إلى أن سبب النزول يفتح باب الفهم ولا يحبس الحكم على صورته الأولى؛ فكل من تعمد أن يخرج في حق الله والخلق أردأ ما عنده مع إمكانه أن يخرج خيراً منه، فقد دخل في المعنى من وجه.

\section{ولستم بآخذيه إلا أن تغمضوا فيه}
استعرض \rawi{الطبري} أقوال السلف في معنى هذه الآية:
\begin{itemize}
    \item فقيل: لا تأخذونه من غرمائكم في حقوقكم إلا مع تنقيص وتغاضٍ.
    \item وقيل: لا تشترونه من السوق بسعر الجيد إلا مع هضم من ثمنه.
    \item وقيل: لا تقبلونه لو أهدي إليكم إلا على استحياء من صاحبه.
\end{itemize}

وهذه الوجوه كلها متقاربة، يجمعها أصل واحد: أن الإنسان لا يرضى لنفسه من حقوقه ما يرضى أن يدفعه في حق الله أو في حقوق الناس إذا كان رديئاً ناقصاً. فكيف يليق به أن يجعل لله أو للمساكين ما يأنف هو من قبوله إلا بغض طرف أو مجاملة أو نقصان؟

\begin{fawaid}[كما لا ترضى الرديء في حقك فلا ترضه في حق الله والخلق]
من أعدل الموازين في أبواب المعاملة والصدقة أن يعرض العبد ما يخرجه على نفسه: هل لو أعطيته أنا على وجه الاستحقاق قبلته راضياً؟ فإن كان الجواب لا، فليتهم إخراجه قبل أن يتهم الآخذ.
\end{fawaid}

\section{الفرق بين الزكاة الواجبة والصدقة التطوعية}
من أنفع ما حرره \rawi{الطبري} هنا تفريقه بين البابين:
\begin{itemize}
    \item فالزكاة الواجبة: لا يجوز أن يخرج فيها الرديء ويمنع أصحاب السهام حقهم من الجيد، لأنهم شركاء في المال الواجب بقدر ما فرض الله لهم.
    \item وأما التطوع: فمع أن إخراج الأجود أحب وأكمل، فإنه لا يحرم إخراج ما دونه إذا كان نافعاً مقبولاً، بل قد يكون أنفع للفقير لكثرته أو لكونه أحسن موقعاً عنده.
\end{itemize}

وهذا التفريق من أجمل ما في الموضع، لأنه يجمع بين تعظيم حق الله في الواجب، وبين فقه المصالح في باب التطوع.

\begin{fawaid}[باب الواجب أضيق من باب التطوع في جنس المخرَج]
ليس حكم الزكاة المفروضة كحكم الصدقة المستحبة؛ فالواجب يجب أن يصان عن التلاعب والتخلص بالرديء، وأما التطوع فينظر فيه أيضاً إلى النفع والمصلحة ما دام المخرج مباحاً نافعاً.
\end{fawaid}

\section{واعلموا أن الله غني حميد}
ختم \rawi{الطبري} هذا الموضع بأن الله غني عن صدقات العباد كلها، وإنما أمرهم بها رحمةً بهم، وتقويةً لضعفائهم، وتكثيراً لأجورهم. فهو سبحانه لا تنفعه طاعة المطيعين، ولا تضره معصية العاصين، لكنه يحمد بما أولى عباده من النعم، وما دلهم عليه من أبواب البر والرحمة.

وهذا من أجمل ما يختم به باب الصدقة: أن العبد إنما يزكي نفسه، ويعالج شحه، ويقدم لآخرته، أما الله سبحانه فغني حميد قبل عطائه وبعده.

\begin{fawaid}[الله غني عن الصدقة وإنما نفعها راجع إلى العبد]
إذا استقر هذا المعنى في القلب سهل الإخلاص، وذهب العجب، وعلم المنفق أنه المنتفع أولاً وآخراً، وأن الله إنما فتح له باب الصدقة رحمة به لا حاجة منه إليه.
\end{fawaid}

\separator

خلاصة هذا اللقاء أن آيات الإنفاق هنا تعلم المؤمن كيف ينفق، ولماذا ينفق، وما الذي يفسد نفقته، وما الذي يقبل منها، وتضرب له من الأمثال ما يجعله يكره الرياء والمنّ والأذى كما يكره ضياع عمره كله عند شدة حاجته إليه. وفيها كذلك تحرير نفيس للفرق بين الزكاة الواجبة والصدقة التطوعية، وربط جميل بين البيان الإلهي للأحكام وبين حقيقة التقوى والورع.
